\item En la producción de rayos X, los electrones son acelerados a través de un alto voltaje y luego se desaceleran al golpear un objetivo.
\begin{enumerate}
    \item Para hacer posible la producción de un rayo X de longitud de onda $\lambda$, ¿cuál es la mínima diferencia de potencial $\Delta V$ requerida?
    \item Explique si esta relación se aplica a otros tipos de radiación electromagnética además de los rayos X.
    \item ¿A qué valor tiende la diferencia de potencial conforme $\lambda$ tiende a cero y
    \item conforme $\lambda$ aumenta sin límite?
\end{enumerate}